\section{Formal properties}
\label{section-formal-cat-cat}

\noindent
In this section we discuss some formal properties of the
$2$-category of categories. This will lead us to the definition
of a (strict) $2$-category later.

\medskip\noindent
Let us denote $\Ob(\textit{Cat})$ the class of all categories.
For every pair of categories
$\mathcal{A}, \mathcal{B} \in \Ob(\textit{Cat})$
we have the ``small'' category of functors
$\text{Fun}(\mathcal{A}, \mathcal{B})$.
Composition of transformation of functors such as
$$
\xymatrix{
\mathcal{A}
\rruppertwocell^{F''}{t'}
\ar[rr]_(.3){F'}
\rrlowertwocell_F{t}
& &
\mathcal{B}
}
\text{ composes to }
\xymatrix{
\mathcal{A}
\rrtwocell^{F''}_F{\ \ t \circ t'}
& &
\mathcal{B}
}
$$
is called {\it vertical} composition. We will use the usual
symbol $\circ$ for this. Next, we will define {\it horizontal}
composition. In order to do this we explain a bit more
of the structure at hand.

\medskip\noindent
Namely for every triple
of categories $\mathcal{A}$, $\mathcal{B}$, and $\mathcal{C}$
there is a composition law
$$
\circ : \Ob(\text{Fun}(\mathcal{B}, \mathcal{C}))
\times
\Ob(\text{Fun}(\mathcal{A}, \mathcal{B}))
\longrightarrow
\Ob(\text{Fun}(\mathcal{A}, \mathcal{C}))
$$
coming from composition of functors. This composition law
is associative, and identity functors act as units. In other
words -- forgetting about transformations of functors --
we see that $\textit{Cat}$ forms a category. How does
this structure interact with the morphisms between functors?

\medskip\noindent
Well, given $t : F \to F'$ a transformation of
functors $F, F' : \mathcal{A} \to \mathcal{B}$ and
a functor
$G : \mathcal{B} \to \mathcal{C}$ we can define
a transformation of functors
$G\circ F \to G \circ F'$. We will denote this
transformation ${}_Gt$. It is given by the formula
$({}_Gt)_x = G(t_x) : G(F(x)) \to G(F'(x))$
for all $x \in \mathcal{A}$.
In this way composition
with $G$ becomes a functor
$$
\text{Fun}(\mathcal{A}, \mathcal{B})
\longrightarrow
\text{Fun}(\mathcal{A}, \mathcal{C}).
$$
To see this you just have to check that
${}_G(\text{id}_F) = \text{id}_{G \circ F}$ and that
${}_G(t_1 \circ t_2) = {}_Gt_1 \circ {}_Gt_2$.
Of course we also have that ${}_{\text{id}_\mathcal{B}}t = t$.

\medskip\noindent
Similarly, given $s : G \to G'$ a transformation of
functors $G, G' : \mathcal{B} \to \mathcal{C}$ and
$F : \mathcal{A} \to \mathcal{B}$ a functor we can define
$s_F$ to be the transformation of functors
$G\circ F \to G' \circ F$ given by
$(s_F)_x = s_{F(x)} : G(F(x)) \to G'(F(x))$
for all $x \in \mathcal{A}$. In this way
composition with $F$ becomes a functor
$$
\text{Fun}(\mathcal{B}, \mathcal{C})
\longrightarrow
\text{Fun}(\mathcal{A}, \mathcal{C}).
$$
To see this you just have to check that
$(\text{id}_G)_F = \text{id}_{G\circ F}$ and that
$(s_1 \circ s_2)_F = s_{1, F} \circ s_{2, F}$.
Of course we also have that $s_{\text{id}_\mathcal{B}} = s$.

\medskip\noindent
These constructions satisfy the additional properties
$$
{}_{G_1}({}_{G_2}t) = {}_{G_1\circ G_2}t,
\ (s_{F_1})_{F_2} = s_{F_1 \circ F_2},
\text{ and }{}_H(s_F) = ({}_Hs)_F
$$
whenever these make sense.
Finally, given functors $F, F' : \mathcal{A} \to \mathcal{B}$,
and $G, G' : \mathcal{B} \to \mathcal{C}$ and transformations
$t : F \to F'$, and $s : G \to G'$ the following
diagram is commutative
$$
\xymatrix{
G \circ F \ar[r]^{{}_Gt} \ar[d]_{s_F}
&
G \circ F' \ar[d]^{s_{F'}} \\
G' \circ F \ar[r]_{{}_{G'}t}
&
G' \circ F'
}
$$
in other words ${}_{G'}t \circ s_F = s_{F'}\circ {}_Gt$.
To prove this we just consider what happens on
any object $x \in \Ob(\mathcal{A})$:
$$
\xymatrix{
G(F(x)) \ar[r]^{G(t_x)} \ar[d]_{s_{F(x)}}
&
G(F'(x)) \ar[d]^{s_{F'(x)}} \\
G'(F(x)) \ar[r]_{G'(t_x)}
&
G'(F'(x))
}
$$
which is commutative because $s$ is a transformation
of functors. This compatibility relation allows us
to define horizontal composition.

\begin{definition}
\label{definition-horizontal-composition}
Given a diagram as in the left hand side of:
$$
\xymatrix{
\mathcal{A}
\rtwocell^F_{F'}{t}
&
\mathcal{B}
\rtwocell^G_{G'}{s}
&
\mathcal{C}
}
\text{ gives }
\xymatrix{
\mathcal{A}
\rrtwocell^{G \circ F} _{G' \circ F'}{\ \ s \star t}
& &
\mathcal{C}
}
$$
we define the {\it horizontal} composition $s \star t$ to be the
transformation of functors ${}_{G'}t \circ s_F = s_{F'}\circ {}_Gt$.
\end{definition}

\noindent
Now we see that we may recover our previously constructed
transformations ${}_Gt$ and $s_F$ as
$ {}_Gt = \text{id}_G \star t $ and $ s_F = s \star \text{id}_F $.
Furthermore, all of the rules we found above are consequences of
the properties stated in the lemma that follows.

\begin{lemma}
\label{lemma-properties-2-cat-cats}
The horizontal and vertical compositions have the following
properties
\begin{enumerate}
\item $\circ$ and $\star$ are associative,
\item the identity transformations $\text{id}_F$
are units for $\circ$,
\item the identity transformations of the identity functors
$\text{id}_{\text{id}_\mathcal{A}}$
are units for $\star$ and $\circ$, and
\item given a diagram
$$
\xymatrix{
\mathcal{A}
\rruppertwocell^F{t}
\ar[rr]_(.3){F'}
\rrlowertwocell_{F''}{t'}
& &
\mathcal{B}
\rruppertwocell^G{s}
\ar[rr]_(.3){G'}
\rrlowertwocell_{G''}{s'}
& &
\mathcal{C}
}
$$
we have $ (s' \circ s) \star (t' \circ t) = (s' \star t') \circ (s \star t)$.
\end{enumerate}
\end{lemma}

\begin{proof}
The last statement turns using our previous notation into the following
equation
$$
s'_{F''}
\circ
{}_{G'}t'
\circ
s_{F'}
\circ
{}_Gt
=
(s' \circ s)_{F''}
\circ
{}_G(t' \circ t).
$$
According to our result above applied to the middle composition
we may rewrite the left hand side as
$
s'_{F''}
\circ
s_{F''}
\circ
{}_Gt'
\circ
{}_Gt
$
which is easily shown to be equal to the right hand side.
\end{proof}

\noindent
Another way of formulating condition (4) of the lemma is
that composition of functors and horizontal composition
of transformation of functors gives rise to a functor
$$
(\circ, \star) :
\text{Fun}(\mathcal{B}, \mathcal{C})
\times
\text{Fun}(\mathcal{A}, \mathcal{B})
\longrightarrow
\text{Fun}(\mathcal{A}, \mathcal{C})
$$
whose source is the product category,
see Definition \ref{definition-product-category}.

